
\documentclass[11pt,a4paper]{article}

% Cargar configuración desde archivo externo
% =========================
% PAQUETES
% =========================
\usepackage[utf8]{inputenc}
\usepackage[T1]{fontenc}
\usepackage[spanish]{babel}
\usepackage{geometry}
\usepackage{graphicx}
\usepackage{xcolor}
\usepackage{tikz}
\usepackage{tabularx}
\usepackage{enumitem}
\usepackage{hyperref}
\usepackage{fancyhdr}
\usepackage{parskip}
\usepackage{ragged2e}
\usepackage{helvet}
\usepackage{titlesec}  % <-- Agregar este paquete
\renewcommand{\familydefault}{\sfdefault}

% Justificacion uniforme y mejor ajuste de lineas para evitar desbordes.
\AtBeginDocument{\justifying}
\setlength{\emergencystretch}{1.5em}

% =========================
% GEOMETRÍA
% =========================
\geometry{
    top=2cm,
    bottom=2.5cm,
    left=2.5cm,
    right=2.5cm
}

% =========================
% COLORES PERSONALIZADOS
% =========================
\definecolor{ieeeOrange}{HTML}{E07830}
\definecolor{ieeeGreen}{HTML}{8FBC8F}
\definecolor{ieeeLightBlue}{HTML}{87CEEB}
\definecolor{ieeeGray}{HTML}{C0C0C0}
\definecolor{titleBrown}{HTML}{5D4E37}
\definecolor{sectionTeal}{HTML}{008080}
\definecolor{linkTeal}{HTML}{007070}
\definecolor{textBrown}{HTML}{716650}

% =========================
% FORMATEO DE SECCIONES
% =========================
\titleformat{\section}
    {\titlerule
    \Large\bfseries\color{ieeeOrange}}  % Formato: grande, negrita, naranja
    {}                                    % Sin número
    {0pt}                                 % Sin espacio entre número y título
    {}                                    % Sin prefijo
\titlespacing*{\section}{0pt}{0.8cm}{0.4cm}  % Espaciado: izq, antes, después

\titleformat{\subsection}
    {\large\bfseries\color{ieeeOrange}}
    {}
    {0pt}
    {}
\titlespacing*{\subsection}{0pt}{0.6cm}{0.3cm}

% =========================
% COMANDOS PERSONALIZADOS
% =========================

% Barra de colores superior
\newcommand{\colorbar}{%
    \noindent
    \begin{tikzpicture}[remember picture, overlay]
        \fill[ieeeOrange] (current page.north west) rectangle ++(5.25cm,-0.4cm);
        \fill[ieeeGreen] ([xshift=5.25cm]current page.north west) rectangle ++(5.25cm,-0.4cm);
        \fill[ieeeLightBlue] ([xshift=10.5cm]current page.north west) rectangle ++(5.25cm,-0.4cm);
        \fill[ieeeGray] ([xshift=15.75cm]current page.north west) rectangle ++(5.25cm,-0.4cm);
    \end{tikzpicture}
}

% Línea de acción (color naranja)
\newcommand{\lineaAccion}[2]{%
    \subsection*{Línea de Acción #1: #2}
}

% Fila de tabla descriptiva
\newcommand{\filaDescriptiva}[2]{%
    \textbf{#1} & #2 \\[0.4cm]
}

\setlist[enumerate,1]{label=\Roman*., font=\bfseries, leftmargin=1.5cm}
\setlist[enumerate,2]{label=\roman*., leftmargin=1cm}

\newcommand{\logos}{
% Logos (placeholder - reemplazar con logos reales)
\noindent
\begin{minipage}[c]{0.2\textwidth}
    \centering
    % Logo IEEE (diamante)
    \includegraphics[width=0.9\linewidth]{src/IEEE.png}
\end{minipage}%
\hfill
\begin{minipage}[c]{0.4\textwidth}
    \centering
    % Logo IEEE Aguascalientes Section
    \includegraphics[width=0.9\linewidth]{src/IEEE_AGS.png}
\end{minipage}%
\hfill
\begin{minipage}[c]{0.25\textwidth}
    \centering
    % Logo Young Professionals
    \includegraphics[width=1.2\linewidth]{src/UAA2.png}

\end{minipage}



}

% =========================
% DOCUMENTO
% =========================
\begin{document}
\color{textBrown}

% ============================================================
% PORTADA
% ============================================================
\colorbar
\pagestyle{empty}

\vspace{1.5cm}

% Logos (placeholder - reemplazar con logos reales)
\noindent
\begin{minipage}[c]{0.2\textwidth}
    \centering
    % Logo IEEE (diamante)
    \includegraphics[width=0.9\linewidth]{src/IEEE.png}
\end{minipage}%
\hfill
\begin{minipage}[c]{0.4\textwidth}
    \centering
    % Logo IEEE Aguascalientes Section
    \includegraphics[width=0.9\linewidth]{src/IEEE_AGS.png}
\end{minipage}%
\hfill
\begin{minipage}[c]{0.25\textwidth}
    \centering
    % Logo Young Professionals
    \includegraphics[width=1.2\linewidth]{src/UAA2.png}

\end{minipage}

\vspace{3cm}

% Título principal
{\fontsize{36}{42}\selectfont\bfseries\color{titleBrown} Plan estratégico de trabajo 2026}

\vspace{0.8cm}

% Subtítulo
{\large Rama Estudiantil UAA Campus Central}

\vspace{0.3cm}

% Línea decorativa
\noindent\rule{7.5cm}{0.5pt}

\vfill

% Información inferior
{\color{sectionTeal}\bfseries IEEE Sección Aguascalientes}\\[0.2cm]
Última fecha de revisión: 6 de octubre 2026\\[0.1cm]
Fecha de autorización:

\newpage

% ============================================================
% DATOS DE IDENTIFICACIÓN
% ============================================================
\pagestyle{fancy}
\fancyhf{}
\renewcommand{\headrulewidth}{0pt}
\fancyfoot[C]{\thepage}

\vspace{0.5cm}

\section{Datos de identificación del documento}

\vspace{0.8cm}

\hyphenpenalty=10000
\exhyphenpenalty=10000

\begin{center}
\begin{tabular}{@{}p{4.5cm} p{8cm}@{}}
    \small Categoría & \small Detalle \\[0.5cm]
    
    \textbf{Título del Plan} & Fortalecimiento Comunitario y Sostenibilidad Operativa 2026 \\[0.5cm]
    
    \textbf{Sección IEEE} & Aguascalientes (Región 9) \\[0.5cm]
    
    \textbf{Coordinación Responsable} & Rama estudiantil IEEE campus UAA \\[0.5cm]
    
    \textbf{Líder del Proyecto} & Sara Fernández Muñoz \\[0.5cm]
    
    \textbf{Fecha de Redacción} & Enero de 2026 \\[0.5cm]
    
    \textbf{Periodo de Aplicación} & Febrero - Febrero 2027 \\[0.5cm]
\end{tabular}
\end{center}


% ============================================================
% CONTENIDO
% ============================================================
\pagestyle{fancy}
\fancyhf{}
\renewcommand{\headrulewidth}{0pt}
\fancyfoot[C]{\thepage}

\vspace{0.5cm}

\section{Introducción y alcance estratégico}
Este plan no solo define las actividades operativas de 2026, sino que también establece una estructura metodológica y de gestión estandarizada para ser fácilmente replicada y transferida a futuros liderazgos. Además, la ejecución de las acciones se detalla a continuación:

Descripción del Evento de Reclutamiento
Como eje central de la estrategia de expansión, se llevará a cabo una jornada de reclutamiento masivo dirigida a las áreas de ingeniería y tecnología.

Sede Propuesta: Auditorio Pedro de Alba, Campus Central UAA.

Público Objetivo: Estudiantes de las carreras de Ingeniería en Sistemas Computacionales, Ingeniería Electrónica e Ingeniería Industrial (seleccionadas por su afinidad con los grupos de afinidad y sociedades técnicas del IEEE).

Dinámica de Integración: El evento consistirá en una presentación de alto impacto sobre los beneficios de la membresía global, el acceso a recursos técnicos y la red de contactos internacional.

Estrategia de Retención y Networking
Para fomentar el sentido de pertenencia desde el primer contacto, el evento concluirá con un espacio de convivencia diseñado para el intercambio de ideas entre aspirantes y miembros actuales:

Aperitivo Institucional: Se ofrecerá café y pan dulce a todos los asistentes, transformando el cierre del evento en un espacio de networking relajado.

Identidad de Marca (Merchandising): Se realizará la entrega de artículos oficiales (merch) de la Rama Estudiantil IEEE a los nuevos interesados, con el fin de proyectar una imagen sólida y atractiva frente a la comunidad universitaria.
\section{Líneas de acción y actividades}

La estrategia operativa se centra en tres líneas de acción fundamentales para el fortalecimiento de la comunidad.

\begin{itemize}[leftmargin=1.5cm]
    \item \textbf{Nota Estratégica:} La ejecución exitosa de estas actividades tiene un doble propósito: la autofinanciación de eventos (ej. Posada) y el aumento de la base de miembros para incrementar la asignación de fondos por parte de la IEEE a la Sección Aguascalientes.
\end{itemize}

\section{Posibles Ideas}


\lineaAccion{I}{Reclutamiento y Difusión}
\renewcommand{\arraystretch}{1.5}
\begin{tabularx}{\textwidth}{@{}p{3cm}X@{}}
    \filaDescriptiva{Evento}{Sesiones de difusión (1ra Edición)-}
    
    \filaDescriptiva{Objetivo Principal}{\textbf{Promover la rama estudiantil,} se busca promover la rama estudiantil en los grupos de ingeniería en electrónica para conseguir una buena base de alumnos a largo plazo.}
    
    \filaDescriptiva{Plan De Acción}{\textbf{Se pretende hacer un} listado de las carreras viables. Así mismo, ubicar profesores o compañeros conocidos de las mismas para facilitar esas reuniones.\newline\newline Modificar la presentación y perfeccionar.}
    
    \filaDescriptiva{Antecedente}{Previamente se había logrado crear un comité para la organización del congreso, entonces se busca retomar esos contactos para facilitar la organización de estas sesiones.}
    
    \filaDescriptiva{Puntos de Mejora}{\textbf{Garantizar la puntualidad} estricta de las empresas.\newline\textbf{Solicitar un \textit{background} más específico} sobre los participantes.}
\end{tabularx}
\renewcommand{\arraystretch}{1.0}
\vspace{1cm}

\lineaAccion{II}{Plan de Remodelación de laboratorios de electrónica.}
\renewcommand{\arraystretch}{1.5}

\begin{tabularx}{\textwidth}{@{}p{3cm}X@{}}
    \filaDescriptiva{Evento}{Diseñar un plan de remodelación para los laboratorios de electrónica de la universidad autonóma.}
    
    \filaDescriptiva{Objetivo Principal}{\textbf{Implementar una mejor organización} en los distintos laboratorios de electrónica de la universidad autonóma.}
    
    \filaDescriptiva{Plan De Acción}{\textbf{Diseñar un plan de trabajo} con recursos de la universidad, enfocándonos en organización y gestión del laboratorio, entendiendo las necesidades y limitaciones para diseñar un plan de trabajo efectivo.}
    
    \filaDescriptiva{Antecedente}{Previamente los laboratorios han sido organizados por alumnos de servicio social implementando pequeñas mejoras, sin embargo, no se ha logrado mantener un orden o una organización a largo plazo.}
    
    \filaDescriptiva{Puntos de Mejora}{\textbf{Establecer un comité de mantenimiento} con alumnos voluntarios.\newline\textbf{Crear un sistema de registro} para el uso y mantenimiento de los laboratorios.}
\end{tabularx}
\renewcommand{\arraystretch}{1.0}

\vspace{1cm}

\lineaAccion{III}{Networking}
\renewcommand{\arraystretch}{1.5}
\begin{tabularx}{\textwidth}{@{}p{3cm}X@{}}
    \filaDescriptiva{Evento}{Encuentros profesionales y sociales}
    
    \filaDescriptiva{Objetivo Principal}{\textbf{Fortalecer la comunidad IEEE} mediante eventos que fomenten la interacción entre miembros, egresados y profesionales de la industria.}
    
    \filaDescriptiva{Plan De Acción}{Organizar al menos un evento de networking por semestre.\newline\newline Invitar a egresados destacados como mentores.}
    
    \filaDescriptiva{Antecedente}{La comunidad ha expresado interés en más oportunidades de conexión profesional.}
    
    \filaDescriptiva{Puntos de Mejora}{\textbf{Establecer seguimiento} posterior a los eventos.\newline\textbf{Crear una base de datos} de contactos profesionales.}
\end{tabularx}
\renewcommand{\arraystretch}{1.0}


\lineaAccion{IV}{Talleres de desarrollo profesional}
\renewcommand{\arraystretch}{1.5}
\begin{tabularx}{\textwidth}{@{}p{3cm}X@{}}
    \filaDescriptiva{Evento}{Talleres que se llevarán a cabo durante el congreso realizado por el centro de Ciencias Básicas (CCB), con temáticas de interés para los miembros.}
    
    \filaDescriptiva{Objetivo Principal}{\textbf{Brindar oportunidades de desarrollo profesional} a los miembros a través de talleres prácticos y relevantes para su formación, Solventar necesidades o carencias identificadas en la comunidad.}
    
    \filaDescriptiva{Plan De Acción}{\textbf{Identificar temáticas de interés} mediante encuestas a los miembros.\newline\newline \textbf{Coordinar con alumnos} para impartir los talleres durante el congreso del CCB.}
    
    \filaDescriptiva{Antecedente}{Alumnos han presentado problemas para ciertas materias, por lo que se busca brindar apoyo adicional a través de talleres.}
    
    \filaDescriptiva{Puntos de Mejora}{\textbf{Establecer seguimiento} posterior a los talleres, buscando crear una comunidad de conocimiento.}
\end{tabularx}
\renewcommand{\arraystretch}{1.0}



\section{Indicadores clave de desempeño (métricas)}

\small
\renewcommand{\arraystretch}{1.5}
\begin{tabularx}{\textwidth}{@{}p{3cm} X X@{}}
    Métrica & Métrica de Medición & Meta de Referencia \\[0.4cm]
    \hline
    \\[-0.3cm]
    
    \textbf{Crecimiento Neto de Miembros} & 
    Número total de nuevos miembros afiliados a la Sección al cierre de 2026. & 
    Incremento mínimo del \textbf{15\%} en la base de miembros activos. \\[0.8cm]
    
    \textbf{Tasa de Conversión (Evento Bienvenida)} & 
    Porcentaje de asistentes al evento de bienvenida que adquieren o renuevan su membresía en los 30 días posteriores. & 
    Tasa de conversión superior al \textbf{35\%} de los asistentes. \\[0.8cm]
    
    \textbf{Generación de Recursos} & 
    Monto total de ingresos brutos generados por la venta de accesos. & 
    Recaudar la totalidad de los fondos necesarios para autofinanciar la \textbf{Posada Anual}. \\[0.8cm]
    
    \textbf{Asistencia \textit{Chill \& Chamba}} & 
    Conteo de asistentes validados. & 
    Superar los \textbf{100 asistentes}. \\[0.8cm]
    
    \textbf{Satisfacción General del Evento} & 
    Promedio de la escala de satisfacción (1 a 5) de las \textbf{Encuestas de Retroalimentación} post-evento. & 
    Puntuación promedio mínima de \textbf{4.5} para todos los eventos. \\[0.8cm]
    
    \textbf{Eficacia del \textit{Networking} (\textit{Chill \& Chamba})} & 
    Porcentaje de encuestados que reportan haber logrado al menos un \textbf{contacto profesional/laboral} relevante. & 
    Lograr que el \textbf{60\%} de los asistentes reporten un contacto valioso. \\[0.5cm]
\end{tabularx}
\renewcommand{\arraystretch}{1.0}
\normalsize


\section{Organizadores y acta de acuerdo}
Este Plan fue redactado por la Ing. Tania Vanessa Ramírez Aguilar, Chair de Young Professionals y Coordinadora de Actividades Empresariales, basándose en los acuerdos alcanzados durante la reunión de organización sostenida el Sábado 4 de octubre de 2025.

La discusión, aprobación del objetivo central y desarrollo de actividades fue consensuada por los siguientes participantes:
Ing. Tania Vanessa Ramírez Aguilar
Ana Sofía Medina Godinez, Tesorera de la Rama Estudiantil del ITA
Ing. David Vitela


\end{document}
